% !TeX root = ../navier-stokes-manuscript.tex
% ============================================================
% 2.1 Vortex Stretching
% ============================================================

The momentum equation has four pieces:
\[
  \underbrace{\partial_t u}_{\text{local acceleration}}
  +
  \underbrace{(u\cdot\nabla)u}_{\text{nonlinear transport}}
  =
  \underbrace{-\nabla p}_{\text{pressure}}
  +
  \underbrace{\nu\Delta u}_{\text{viscous diffusion}}.
\]
The central regularity problem lies in the competition between nonlinear transport and viscosity.
The hope is simple: viscosity smooths the fluid faster than nonlinear transport can sharpen it.
If that remains true at every scale, the velocity stays smooth.

\subsection{Vortex Stretching}\label{sub:ThePerfectlyStretchingVortex}

The surrounding flow pulls a narrow vortex outward along its axis.
On a short material segment where the strain is spatially uniform, let \(e(t)\) be its instantaneous unit tangent, aligned with a positive strain direction:
\[
  S:=\frac12\bigl(\nabla u+(\nabla u)^{\mathsf T}\bigr),
  \qquad
  Se(t)=\sigma(t)e(t),
  \qquad
  \sigma(t)>0.
\]
The material-line law gives its length \(L(t)\) by
\[
  \frac{\dot L(t)}{L(t)}
  =
  e(t)\cdot Se(t)
  =
  \sigma(t).
\]
The vortex lengthens.
But it is the same incompressible fluid.
Its material volume \(\mathcal V\) cannot change.
The effective cross-sectional area \(A(t):=\mathcal V/L(t)\) and effective radius \(r_{\mathrm{eff}}(t):=\sqrt{A(t)/\pi}\) then satisfy
\[
  \nabla\cdot u=0,
  \qquad
  \frac{d}{dt}(A L)=0,
  \qquad
  \frac{\dot A}{A}=-\sigma(t),
  \qquad
  \frac{\dot r_{\mathrm{eff}}}{r_{\mathrm{eff}}}
  =
  -\frac{\sigma(t)}{2}.
\]
As the length grows, the effective cross-section falls.
The vortex becomes narrower.

The same axial strain acts on the rotation carried by the vortex.
With \(\omega:=\nabla\times u\),
\[
  \frac{D\omega}{Dt}
  =
  S\omega+\nu\Delta\omega,
  \qquad
  \frac{D}{Dt}
  :=
  \partial_t+u\cdot\nabla.
\]
When the vorticity along the tracked trajectory remains aligned with the tangent, \(\omega=|\omega|e(t)\),
\[
  S\omega
  =
  \sigma(t)\omega,
  \qquad
  \frac12\frac{D|\omega|^2}{Dt}
  =
  \sigma(t)|\omega|^2
  +
  \nu\,\omega\cdot\Delta\omega.
\]
Stretching contributes \(\sigma(t)|\omega|^2>0\).
Viscosity has no fixed pointwise sign; to stop the amplification, it must offset that contribution.
When their sum remains positive along the material trajectory, \(|\omega|\) grows while \(r_{\mathrm{eff}}\) falls.
At the same time, \Cref{prop:ClassicalKineticEnergyIdentity} keeps the total kinetic energy finite:
\[
\begin{aligned}
  \norm{u(t)}_2
  &\leq
  \norm{u_0}_2
  <
  \infty,
  \\
  \frac{|\omega|}{\sqrt2}
  &=
  \left|\frac12\bigl(\nabla u-(\nabla u)^{\mathsf T}\bigr)\right|_{\mathrm F}
  \leq
  |\nabla u|_{\mathrm F},
  \\
  \frac{D}{Dt}
  \left|\frac12\bigl(\nabla u-(\nabla u)^{\mathsf T}\bigr)\right|_{\mathrm F}
  &>
  0.
\end{aligned}
\]
The rotational gradients rise.

\divider{\NavierStokes{} Scaling}

The radius is falling while stretching and viscosity are still competing over the same rotation.
At a time \(t_0<T_*\), write \(\ell:=r_{\mathrm{eff}}(t_0)\).
For \(\lambda>1\), the exact \NavierStokes{} comparison at width \(\lambda^{-1}\ell\) is
\[
  u_\lambda(x,t)
  :=
  \lambda u(\lambda x,\lambda^2t),
  \qquad
  p_\lambda(x,t)
  :=
  \lambda^2p(\lambda x,\lambda^2t).
\]
At time \(\lambda^{-2}t_0\), the corresponding element in \(u_\lambda\) has effective radius \(\lambda^{-1}\ell\).
This is another solution at another scale, not a later state of the original vortex.
Each momentum term transforms with it:
\[
\begin{aligned}
  \partial_tu_\lambda(x,t)
  &=
  \lambda^3(\partial_tu)(\lambda x,\lambda^2t),
  \\
  (u_\lambda\cdot\nabla)u_\lambda(x,t)
  &=
  \lambda^3\bigl((u\cdot\nabla)u\bigr)(\lambda x,\lambda^2t),
  \\
  \nabla p_\lambda(x,t)
  &=
  \lambda^3(\nabla p)(\lambda x,\lambda^2t),
  \\
  \nu\Delta u_\lambda(x,t)
  &=
  \lambda^3\nu(\Delta u)(\lambda x,\lambda^2t).
\end{aligned}
\]
Every term gains the same \(\lambda^3\).
At the finer width viscosity is stronger, but nonlinear sharpening is stronger by exactly the same factor.
The comparison has given viscosity no relative advantage.

\divider{Critical Height}

That tie leaves one Sobolev height untouched.
With \(\Lambda:=(-\Delta)^{1/2}\),
\[
  \widehat{u_\lambda}(\xi,t)
  =
  \lambda^{-2}\widehat u(\lambda^{-1}\xi,\lambda^2t),
\]
and a change of variables gives
\[
  \norm{u_\lambda(t)}_{\dot H^s}^2
  =
  \lambda^{2s-1}
  \norm{u(\lambda^2t)}_{\dot H^s}^2.
\]
At \(s=0\), concentration lowers kinetic energy.
At \(s=1\), it raises the first-derivative norm.
Between them, the exponent vanishes at \(s=\tfrac12\).
Set
\[
  H(t)
  :=
  \frac12\norm{u(t)}_{\dot H^{1/2}}^2
  =
  \frac12\norm{\Lambda^{1/2}u(t)}_2^2.
\]
Writing \(H_\lambda\) for the same quantity evaluated on \(u_\lambda\),
\[
  \norm{u_\lambda(t)}_2^2
  =
  \lambda^{-1}\norm{u(\lambda^2t)}_2^2,
  \qquad
  H_\lambda(t)=H(\lambda^2t),
  \qquad
  \norm{\nabla u_\lambda(t)}_2^2
  =
  \lambda\norm{\nabla u(\lambda^2t)}_2^2.
\]
The comparison can move the structure to a finer width, but it cannot move \(H\).

\divider{Nonlinear Production versus Viscous Dissipation}

If \(H\) rises, the original equation has done it.
Its signed nonlinear contribution is
\[
  \mathcal P_{1/2}[u](t)
  :=
  -\operatorname{Re}
  \pair
    {\Lambda^{1/2}u(t)}
    {\Lambda^{1/2}\bigl((u(t)\cdot\nabla)u(t)\bigr)}_{L^2}.
\]
The pressure pairing vanishes because \(\nabla\cdot u=0\), and the Laplacian contributes \(-\nu\norm{\Lambda^{3/2}u}_2^2\), leaving
\[
  H'(t)
  +
  \nu\norm{\Lambda^{3/2}u(t)}_2^2
  =
  \mathcal P_{1/2}[u](t).
\]
The vortex problem now has an exact form: the height rises only when nonlinear production exceeds viscous dissipation.
Under the same rescaling,
\[
  \mathcal P_{1/2}[u_\lambda](t)
  =
  \lambda^2\mathcal P_{1/2}[u](\lambda^2t),
  \qquad
  \nu\norm{\Lambda^{3/2}u_\lambda(t)}_2^2
  =
  \lambda^2\nu\norm{\Lambda^{3/2}u(\lambda^2t)}_2^2.
\]
Production and dissipation again gain the same power.
Their common \(\lambda^2\) has moved the unresolved contest onto the shortened time coordinate \(t\mapsto\lambda^{-2}t\).

\divider{The Heat Clock}

Viscosity carries its own version of that contracted clock.
For the linear heat equation \(v_t=\nu\Delta v\),
\[
  \widehat v(\xi,t+\delta t)
  =
  e^{-\nu|\xi|^2\delta t}\widehat v(\xi,t).
\]
A scale-localized structure of width \(r\) occupies frequencies \(|\xi|\simeq r^{-1}\).
Its viscous change becomes order one when
\[
  \nu|\xi|^2\delta t\simeq1,
\]
so the linear comparison time is
\[
  \tau_\nu(r)
  :=
  \frac{r^2}{\nu}.
\]
At width \(\lambda^{-1}r\),
\[
  \tau_\nu(\lambda^{-1}r)
  =
  \lambda^{-2}\tau_\nu(r).
\]
The finer comparison has not strengthened viscosity relative to sharpening.
It has shortened the clock for both by the same square.

\divider{The Zeno Cascade}

Repeated halving makes those clocks a geometric sequence.
For
\[
  r_n:=2^{-n}r_0,
  \qquad
  \tau_n:=\frac{r_n^2}{\nu},
\]
the comparison clocks satisfy
\[
  \tau_n
  =
  \frac{r_0^2}{\nu}4^{-n},
  \qquad
  \sum_{n=0}^{\infty}\tau_n
  =
  \frac{4r_0^2}{3\nu}
  <
  \infty.
\]
The clocks are summable.
That arithmetic leaves room for infinitely many finer comparisons before a finite time, but it does not turn them into one fluid history.
For these clocks to describe a candidate singularity, one \NavierStokes{} solution would have to pass through widths
\[
  r_0>r_1>r_2>\cdots\longrightarrow0
\]
at times satisfying the additional schedule
\[
  t_0<t_1<t_2<\cdots\uparrow T_*<\infty,
  \qquad
  t_{n+1}-t_n\asymp\tau_n,
\]
with comparison constants independent of \(n\).
Under that hypothesis, the remaining time is
\[
  T_*-t_n
  \asymp
  \sum_{k=n}^{\infty}\tau_k
  =
  \frac{4r_n^2}{3\nu}.
\]
At the \(n\)-th width, both the next transition and the total time left are of order \(r_n^2/\nu\).
The same velocity field would have to change across intervals that collapse to zero with its spatial scale.
